Il primo metodo utilizzato consiste nel fare un video continuo attraverso la lente grandengolo contenuta nel cellulare per riuscire a visualizzare al meglio il primo ordine di difrazione. Una volta finito i vari processi spiegati nella sezione precedente, estraggo i frame migliori dal video svolto dal sensore durante tutta la durata dell'esperimento ottenendo le due seguenti immagini:
Nella prima immagine è possibile visualizzare la diffrazione prodotta dal primo reticolo olografico (con $d$ noto) e nella seconda quella prodotta dal secondo reticolo olografico (con $d'$ ignoto).